
\documentclass[authoryear,preprint,review,times,10pt]{elsarticle}

\usepackage{amsmath}    
\usepackage{amsthm}    
\usepackage{amssymb}     
\usepackage{amsfonts}  
\usepackage{bm}         

\usepackage{booktabs}   
\usepackage{multirow}    
\usepackage{longtable}   
\usepackage{threeparttable} 
\usepackage{tabularx}    
\usepackage{caption}      
\usepackage{adjustbox}    
\usepackage{pdflscape}    
\usepackage{enumitem}     
\usepackage{lineno}       
\usepackage{setspace}     
\doublespacing           

\usepackage{graphicx}
\usepackage{subfig}
\graphicspath{{simulation/figures/}}
\usepackage{tikz}
\usetikzlibrary{positioning, shapes.geometric, arrows}

\usepackage{geometry}
\geometry{
    left=2.5cm,
    right=2.5cm,
    top=2.5cm,
    bottom=2.5cm,
}

\usepackage{hyperref}    

\makeatletter
\def\ps@pprintTitle{
  \let\@oddhead\@empty
  \let\@evenhead\@empty
  \def\@oddfoot{\hfil\thepage\hfil}
  \let\@evenfoot\@oddfoot}
\def\pprintMaketitle{\clearpage
  \iflongmktitle\if@twocolumn\let\columnwidth=\textwidth\fi\fi
  \resetTitleCounters
  \def\baselinestretch{1}
  \printFirstPageNotes 
  \begin{center}
   \thispagestyle{pprintTitle}
   \def\baselinestretch{1}
    \Large\@title\par\vskip18pt
    \normalsize\elsauthors\par\vskip10pt
    \normalsize\itshape\elsaddress\par\vskip12pt
    \ifvoid\absbox\else\unvbox\absbox\par\vskip10pt\fi
    \ifvoid\keybox\else\unvbox\keybox\par\vskip10pt\fi
  \end{center}
  \gdef\thefootnote{\arabic{footnote}}%
  }
\makeatother

\journal{}

\begin{document}

\newtheorem{prop}{Proposition}

\modulolinenumbers[1]
\runninglinenumbers
\begin{frontmatter}

\title{Biform Game Analysis of Waste Concrete Recycling Management Considering the Cap-and-trade Mechanism}

\author[1]{Zhe Zhang}
\author[2]{Qinghua He, Ph.D.}
\author[3]{Zilun Wang, Ph.D.}

\affiliation[1]{Ph.D. Student; School of Economics and Management, Tongji University, Shanghai 200092, China; Email: 2431220@tongji.edu.cn}
\affiliation[2]{Professor; School of Economics and Management, Tongji University, Shanghai 200092, China; Email: heqinghua@tongji.edu.cn}
\affiliation[3]{Lecturer; School of Civil Engineering, Suzhou University of Science and Technology, Suzhou 215011, China (Corresponding author); Email: zylenw97@usts.edu.cn}


\begin{abstract}

Waste concrete recycling is essential for sustainable construction and resource management. As the main participants, concrete manufacturers and recyclers face critical management challenges, needing to properly balance pricing competition with research and development cooperation. Additionally, cap-and-trade mechanism as the policy tool further complicates their economic decisions. To address these issues, we develop a biform game model integrating cooperative and non-cooperative games under the cap-and-trade mechanism. Compared to a pure non-cooperative game model. Results show that whether enterprises choose to cooperate depends on their evaluation perspective, as cooperation is not always optimal from an economic profit standpoint due to external factors like government subsidies and outsourcing fees. It remains consistently the most beneficial decision for improving waste concrete conversion and reducing carbon emissions from the environment and resource conversion perspective. Furthermore, while the cap-and-trade mechanism raises costs, cooperation effectively buffers against carbon market volatility. This research provides actionable insights for environmental policy design and enterprise management decisions in waste concrete recycling.

\end{abstract}

\begin{keyword}
Waste concrete recycling \sep Biform Game \sep Cap-and-trade Mechanism \sep Cooperation
\end{keyword}

\end{frontmatter}



\section{Introduction}
The efficient recycling and utilization of construction and demolition (C\&D) waste is a vital pathway to achieving resource efficiency and circularity in the construction sector \citep{yuan2025Exploring}. Rapid development of the construction industry in recent decades has led to a sharp increase in construction and demolition (C\&D) waste generation, with waste concrete accounting for more than half of total C\&D waste \citep{zhao2020Use}. Given the environmental and social damage caused by this waste, urgent management of its recycling is required \citep{wu2019Offsite}. Waste concrete recycling has demonstrated potential to alleviate environmental pressures and shortages of construction raw materials due to its regenerative and utilization value \citep{mostert2021Climate}. Many countries, particularly developed nations with limited natural resources, have widely adopted recycled concrete made by the waste concrete \citep{yenipazarli2019Incentives}. However, promotion of waste concrete recycling management faces obstacles including immature technologies, unclear economic profit, and incomplete policies, resulting in low enterprise participation \citep{liu2021Explore,ma2020Challenges}. The waste concrete recycling process involves two major participants: concrete manufacturers and recyclers \citep{he2020Investigation,lu2015Stakeholders}. Confronted with high technical costs and market uncertainties, concrete manufacturers and recyclers often lack the capacity to independently overcome these barriers and require inter-enterprise coordination to achieve complementary advantages, risk sharing, and maximum economic profit from waste concrete recycling \citep{bao2019Procurement}. Therefore, how to effectively coordinate the waste concrete recycling management process has become a critical management challenge constraining the development of the waste concrete recycling industry.

The waste concrete recycling process generally involves collecting raw waste concrete from the market and the treatment \citep{he2020Investigation}. In this process, concrete manufacturers possess the channel advantage. They can collect waste concrete from construction sites using their delivery trucks on return trips, which keeps their transportation costs relatively low \citep{besklubova2026Establishing}. On the other hand, the recyclers use dedicated equipment to process the collected waste into recycled concrete aggregates\citep{chen2025Exploring}. During these interactions, the two participants engage in price competition at the front end to secure limited waste concrete resources \citep{he2020Investigation}. At the same time, they can choose to cooperate at the back end by sharing the research and development investment required to improve these processing technologies \citep{liu2019Expand}. These intertwined behaviors form a complex co-opetition relationship. However, both participants are often hesitant to actually commit to the R\&D investment cooperation \citep{bao2020Developing}. This is because their competitive and cooperative behaviors mutually influence each other. Pricing battles at the front end impact the final profit distribution from R\&D investment sharing at the back end, making the ultimate economic profit of working together unclear \citep{li2020Key}. While previous research has explored recycling strategies, most studies treat collection and treatment or competition and cooperation as separate stages \citep{hu2021Dynamic,yuan2023Differentiated}. They fail to clarify whether these companies should cooperate in such a competitive environment and how they should adjust their prices based on this cooperation to facilitate the recycling of waste concrete resources. Therefore, it is necessary to explore how the interaction between pricing competition and R\&D investment cooperation affects the decisions of concrete manufacturers and recyclers, as well as the impact on the recycling of waste concrete resources.

The environmental impact of the recycling process also complicates the co-opetition relationships. Although waste concrete recycling offers advantages in conserving virgin resources, the treatment process is often accompanied by substantial carbon emissions \citep{difilippo2019Impacts}. Current mainstream waste concrete recycling technologies primarily rely on crushing, separation, and calcination-based thermal treatment, which still generate considerable carbon dioxide \citep{oh2021Proposal}. In this context, technological research and development (R\&D) is designed not only to maximize the conversion rate of waste concrete into high-quality recycled aggregates, but also to mitigate the carbon emissions generated during these processing stages.
To curb industrial carbon emissions, the cap-and-trade mechanism has been widely implemented globally as an effective market-based tool that balances economic profit with environmental protection \citep{oh2021Proposal,yenipazarli2019Incentives}. Waste concrete recycling is also an industry widely affected by the cap-and-trade mechanism.
For the waste concrete recycling industry, carbon quota constraints and fluctuations in trading prices reshape the cost structures and profit margins of participating enterprises, compelling them to consider the economic implications of carbon emissions when making decisions \citep{zheng2025Coopetition}. However, existing research on waste concrete recycling management and economic decision-making has largely overlooked this aspect. The impacts of the cap-and-trade mechanism on these co-opetition relationships remain unclear.
There is a gap in understanding how manufacturers and recyclers should adjust their cooperation strategies under the constraints of environmental policies to ensure the conversion of waste concrete and profitability. Therefore, incorporating the cap-and-trade mechanism into the research framework to investigate how it influences the economic profit of waste concrete recycling and the decision-making behavior of enterprises is necessary.

This study proposes a two-stage biform game model for waste concrete recycling that incorporates both non-cooperative and cooperative strategies to analyze the co-opetition relationships of concrete manufacturers and recyclers. Non-cooperative strategies represent recycling pricing competition between concrete manufacturers and recyclers, which do not yield final equilibrium solutions but rather establish a competitive context for their R\&D investment cooperation decisions. Simultaneously, cooperative strategies are jointly determined by both enterprises. The profit distribution resulting from the cooperative strategies subsequently becomes the critical profit function for resolving the final non-cooperative strategies. Furthermore, we compare this co-opetition strategy with a non-cooperative game model. Specifically, this study addresses the following research questions:

\begin{itemize}
\item \textbf{RQ1:} Considering cap-and-trade mechanism, should concrete manufacturers and recyclers engage in cooperation within the competitive environment?
\item \textbf{RQ2:} What are the impacts of the cap-and-trade mechanism on the decisions of concrete manufacturers and recyclers in waste concrete recycling?
\end{itemize}

The remainder of this paper is organized as follows. Section 2 reviews the relevant literature. Section 3 presents the model formulation, including problem description, assumptions, and profit functions. Section 4 conducts numerical simulations to verify the theoretical results and analyze the effects of key parameters. Section 5 concludes the paper with recommendations for future research.

\section{Literature review}
\subsection{Waste concrete recycling management and economic decision}
In the waste concrete recycling supply chain, existing literature exploring their management and economic decision can be categorized into three main research themes: reverse logistics network design and operation management, government subsidy and incentive mechanism design, and manufacturer-recycler competition and cooperation relationships.

Regarding reverse logistics network design and operation management, researchers primarily focus on optimizing physical flows and capacity. \cite{rahimi2018Sustainable} developed a multi-period mixed integer linear programming (MILP) model to optimize the network design for C\&D waste recycling. Similarly, \cite{xu2019Reverse} proposed a capacity expansion pricing mechanism tailored for recycling networks under demand uncertainty, while \cite{liu2019Expand} identified key cost drivers through a comprehensive network cost analysis. However, these operational models treat the recycling network as a centrally coordinated system, overlooking the decentralized decision-making processes and the profit-driven strategic conflicts between independent concrete manufacturers and recyclers.

Within the theme of government subsidy and incentive mechanism design, studies emphasize utilizing external policy interventions to stimulate market participation. \cite{yuan2023Differentiated} investigated differentiated subsidy schemes under heterogeneous consumer quality perceptions, demonstrating that uniform subsidies often fail to achieve optimal outcomes. \cite{su2024Stakeholder} evaluated the impact of government intervention through quota trading mechanisms on recycling rates, and \cite{zheng2024Optimal} examined how various incentive combinations influence stakeholder participation. 
Although existing research confirms the effectiveness of macro-level policy tools, it often assumes uniform cost structures among supply chain participants. The literature inadequately explains how enterprises' specific operational advantages dictate varying responses to environmental policies.

Finally, concerning the manufacturer-recycler competition and cooperation relationships of the concrete manufacturer and the recycler, scholarship has begun to unpack micro-level strategic behaviors. \cite{he2020Investigation} analyzed optimal decisions under different market scenarios to identify the conditions under which cooperation dominates pure competition. \cite{lin2025Cooperation} revealed the cooperation dilemma (prisoner's dilemma) in C\&D waste recycling, showing that unilateral strategies frequently lead to suboptimal industry outcomes. Furthermore, \cite{li2025Mechanism} employed evolutionary game theory to study the long-term dynamics of manufacturer-recycler relationships. While these contributions provide essential theoretical guidance for the decision-making of concrete manufacturers and recyclers in waste concrete recycling, existing research predominantly focuses on single-scenario analyses or separate scenario comparisons. They fail to capture the simultaneous interaction patterns of competition and cooperation inherent in this process, struggling to resolve the hesitation of R\&D investment cooperation within the industry.

\subsection{Biform games and recycling supply chain coordination}
The theoretical foundation of this study rests on the biform game model, which was originally formalized by \cite{brandenburger2007Biform} to analyze scenarios where firms simultaneously engage in competition and cooperation. This model captures the coopetition relationships by integrating non-cooperative game theory (first stage) with cooperative game theory (second stage). In the first stage, enterprises make strategic moves to shape the competitive environment; in the second stage, they allocate the created surplus through bargaining mechanisms, typically using mechanisms such as the Shapley value.

The application of biform games to supply chain coordination has received substantial attention in recent years. \cite{feess2014Surplus} developed a pioneering supply chain biform game model where firms make non-cooperative investment decisions in the first stage and divide surplus via the Shapley value in the second stage, revealing that all firms' investment incentives fall below the socially optimal level. \cite{fiala2022Modelling} applied biform games to model co-opetition in network industries, distinguishing between sequential and simultaneous game structures. More recently, \cite{zheng2023New} proposed a biform game-based investment incentive mechanism for eco-efficient innovation in supply chains.

Extending this framework to reverse supply chain and recycling contexts, several studies have examined how enterprises coordinate collection and treatment activities under competitive pressures. \cite{jia2023Novel} employed a biform game to coordinate a multi-agent reverse supply chain with fairness concerns, demonstrating that the biform game mechanism could improve all supply chain members' outcomes compared to pure non-cooperative equilibria. \cite{gui2018Design} analyzed design incentives under collective extended producer responsibility using biform games, focusing on how processor capacity sharing affects design decisions and alliance stability in recycling networks. \cite{wang2026Integration} explicitly modeled collection competition and processing cooperation by constructing five game scenarios in a reverse supply chain with horizontal competitive recycling, revealing that full cooperation maximizes total system profit while different allocation methods favor different participants. The integration of biform games with environmental policy mechanisms has also been explored. \cite{zheng2024New} developed a biform game-based coordination mechanism for a carbon complementary supply chain under hybrid carbon regulations, identifying threshold conditions for regulatory intensity beyond which cooperative investment yields Pareto improvements.

Therefore, the biform game framework is well-suited for analyzing waste concrete recycling, where competition in collection markets coexists with cooperation in treatment activities. This analytical approach enables simultaneous examination of how carbon cap-and-trade mechanisms influence collection pricing strategies and how R\&D investment cooperation can be designed to improve individual profits for both concrete manufacturers and recyclers. We develop a biform game model tailored to the waste concrete recycling context to explore their optimal economic decisions.

\subsection{The cap-and-trade mechanism and recycling}
The cap-and-trade mechanism sets a carbon emission cap while allowing enterprises to freely trade emission allowances, incentivizing enterprises to pursue lower-cost emission reduction pathways while providing a flexible balance between economic profit and environmental protection \citep{li2019How}. The implementation of cap-and-trade mechanisms has reshaped the cost structures and operational boundaries of recycling supply chains. A primary stream of literature evaluates the effectiveness of this market-based instrument in balancing economic profitability with environmental compliance. Studies have demonstrated that cap-and-trade policies, when appropriately parameterized, can effectively incentivize emission reductions without severely compromising system performance \citep{bai2022Distributionally}. Furthermore, comparative analyses indicate that the optimal regulatory choice—whether carbon taxes, strict caps, or cap-and-trade—is highly contingent upon market structure, carbon price levels, and firms' technological innovation capacities \citep{lyu2022Manufacturers,huang2024Optimal}. These foundational studies confirm that carbon quotas are no longer mere exogenous constraints, but endogenized variables that dictate recycling participation.

Building on this macro-policy premise, recent scholarship has increasingly focused on how carbon trading influences strategic interactions and mode selection in recycling networks. Under the dual pressures of economic profit and carbon compliance, researchers have highlighted the necessity of joint recycling modes. For instance, dynamic carbon pricing has been shown to induce symbiotic cooperation between manufacturers and recyclers, often outperforming solitary recycling efforts under policy interventions \citep{cao2023Developing,wu2024Electric}. Additionally, optimal recycling channel selection and supply chain pricing strategies are heavily modulated by the allocation of carbon quotas \citep{zhang2022Recycling}.

Despite these valuable insights, a misalignment exists between current carbon-constrained recycling models and the engineering realities of waste concrete. Existing literature predominantly investigates high-value-density, low-transport-friction products (e.g., electronic waste and EV batteries). In contrast, waste concrete is characterized by low unit value, massive physical volume, and highly asymmetrical transportation costs (i.e., the concrete manufacturer's backhaul advantage versus the recycler's dedicated logistics). Furthermore, the processing of waste concrete is inherently carbon-intensive, making the recycler highly sensitive to carbon price volatility. Current models fail to capture how the cap-and-trade mechanism specifically couples with these unique logistical asymmetries to influence front-end pricing competition and back-end R\&D investment. Consequently, extending the cap-and-trade analysis to the waste concrete context recycling is imperative to formulate effective, industry-specific environment and economic strategies.

\section{Model formulation}
\subsection{Problem description}
We consider a two-stage waste concrete recycling supply chain consisting of a traditional concrete manufacturer M and a specialized waste concrete recycler R. Both M and R participate in the collection of waste concrete from waste generators, while recycler R possesses specialized technologies and equipment to process the collected waste into recycled concrete aggregate (RCA). Recycler R generates carbon dioxide emissions during the operations, and therefore must consider the impact of cap-and-trade policies on its economic decisions. In this section, we develop and compare two different game-theoretic models: (1) a pure non-cooperative waste concrete recycling model considering the cap-and-trade mechanism, and (2) a biform game model combining R\&D investment cooperation with the cap-and-trade mechanism.

In the non-cooperative game model, there is no R\&D investment sharing between the concrete manufacturer and the recycler. Recycler R independently bears all the R\&D investment for upgrading the processing technology. The decision sequence is as follows: manufacturer M and recycler R simultaneously determine their respective collection prices to compete for waste concrete resources in the market; subsequently, given these collection prices, recycler R determines the optimal RCA conversion rate to maximize its own profit.

In the biform game model, the interaction involves both pricing competition and R\&D investment cooperation. The game process is modeled in two stages. First, in the non-cooperative stage, concrete manufacturer M and recycler R determine their respective collection prices to form a competitive situation. Second, in the cooperative stage, under any given competitive situation, the two participants jointly form an alliance to determine the optimal conversion rate and the R\&D investment sharing proportion. We calculate the characteristic values of each coalition and apply the Shapley value method to determine the profit distribution values ($\varphi_M$ and $\varphi_R$) for manufacturer M and recycler R, respectively. Finally, using these profit distribution values as the payoff functions, we solve the first-stage non-cooperative game backward to derive the optimal collection prices, conversion rate, and final profit allocations.

The following assumptions are made:

\textbf{Assumption 1.} Manufacturer $M$ and recycler $R$ set collection prices $(p_M, p_R)$ for waste concrete generators. Since both parties participate in collection, there exists competitive pressure in the recycling market. Let $\beta$ denote the competition intensity between $M$ and $R$. Additionally, the collection quantity depends on the collection price, with $\alpha$ representing the sensitivity of waste generators to collection prices. We assume sufficient waste concrete available in the market and $R$ has sufficient capacity to process all collected quantity. Referring to \citep{chen2025Lowcarbon} and \citep{zhang2024Selection}'s research, the collection quantity function follows:
\begin{equation}
q_i = \alpha p_i - \beta p_j \quad (i,j \in \{M, R\}, i \neq j)
\end{equation}
where $\alpha > \beta > 0$ ensures positive collection quantities. A larger $\alpha$ indicates higher sensitivity to prices and thus greater collection quantity, and a larger $\beta$ indicates more intense competition between $M$ and $R$.

\textbf{Assumption 2.} The recycling processing capability of recycler $R$ is characterized by the RCA conversion rate $k \in [0,1]$. Improving this conversion rate (e.g., through enhanced mortar removal technology or multi-stage air separation) requires fixed investment in research and development. Based on the research by \cite{jones2011Information} and \cite{chen2024Risk}, we assume that the R\&D investment cost function is:
\begin{equation}
C(k) = \frac{1}{2} c k^2
\end{equation}
where $c$ is the R\&D investment cost coefficient.

\textbf{Assumption 3.} In practice, traditional concrete manufacturers often lack deep-processing environmental equipment or are constrained by site and environmental qualifications. Therefore, manufacturer $M$ outsources the collected waste concrete to recycler $R$ for centralized processing. Manufacturer $M$ pays a unit outsourcing fee $w$ to recycler $R$, while recycler $R$ incurs actual operational costs $m$ (e.g., utilities, equipment depreciation, labor) per unit of waste concrete processed.

\textbf{Assumption 4.} Assume $\theta$ is the market price per unit of recycled concrete aggregate (RCA) when perfectly converted from waste concrete. Under the actual RCA conversion rate $k$ achieved by recycler $R$, the actual extracted value per unit of waste concrete is $\theta k$. For manufacturer $M$, this value represents the cost saving from substituting natural aggregates with RCA. For recycler $R$, it represents the selling price of processed RCA. To encourage construction waste recycling under the zero-waste city initiative, the government provides a subsidy $s$ per unit of waste concrete collected, which is distributed to the respective collector based on their collection quantity.

\textbf{Assumption 5.} To obtain high-quality recycled aggregate and maintain stable outsourcing cooperation, manufacturer $M$ is willing to share the R\&D investment for upgrading the processing technology with recycler $R$. Let manufacturer $M$ share a proportion $(1-n)$ and recycler $R$ bear a proportion $n$, where $n \in [0,1]$. When $n=1$, all investment costs are borne solely by recycler $R$.

\textbf{Assumption 6.} Both manufacturer $M$ and recycler $R$ incur unit transportation costs for collecting waste concrete from generators\citep{chileshe2016Analysis}. The physical properties of waste concrete (heavy, large volume, low value density) make transportation costs a significant component of total recycling costs. Due to the backhaul logistics advantage in the construction industry, manufacturer $M$'s concrete delivery trucks can collect waste concrete on return trips from construction sites, resulting in a lower transportation cost $\mu_M$. In contrast, recycler $R$ must dispatch dedicated vehicles for collection, incurring a higher transportation cost $\mu_R$. We assume $\mu_M < \mu_R$. This cost asymmetry represents the engineering management characteristic of waste concrete recycling.

\textbf{Assumption 7.} Under the cap-and-trade mechanism, recycler $R$ is subject to government-allocated carbon emission quotas and engages in carbon trading. In practice, when enterprises invest in R\&D for waste concrete processing technologies, they strategically pursue a dual objective: maximizing the resource utilization rate while simultaneously minimizing treatment-related carbon emissions. To capture this synergistic effect of technological upgrading, the actual carbon emissions generated by recycler R processing a unit of waste concrete is established as a dynamic function of the RCA conversion rate $k$, denoted by $e(k)$. Specifically, we assume a linear emission reduction function:
\begin{equation}
e(k) = e_0 - \gamma k
\end{equation}
where $e_0$ is the baseline unit carbon emissions (when $k = 0$) and $\gamma$ is the technical emission reduction coefficient, representing the unit carbon emission reduction achieved per unit increase in the conversion rate $k$. To ensure a non-negative effective emission rate and a feasible interior conversion rate, we impose $0 < k \leq 1$ and $e(k) = e_0 - \gamma k \geq 0$. The total processing quantity is $q_{total} = q_M + q_R = (\alpha - \beta)(p_M + p_R)$. Let $p_c$ denote the trading price per unit of carbon emission in the carbon market. Let $G$ denote the total carbon emission quota allocated by the government to recycler $R$ in a compliance cycle. The net carbon trading cost incurred by recycler $R$ is:
\begin{equation}
p_c \left[ e(k) \cdot q_{total} - G \right] = p_c \left[ (e_0 - \gamma k)(\alpha - \beta)(p_M + p_R) - G \right]
\end{equation}

\begin{table}[!htbp]
\centering
\captionsetup{font=normalsize, labelsep=period}
\caption{Notation and definitions}
\label{tab:notation}
\normalsize
\begin{threeparttable}
\begin{tabular*}{0.9\textwidth}{@{\extracolsep{\fill}}lc}
\toprule
\textbf{Symbol} & \textbf{Definition} \\
\midrule
$\theta$ & Unit potential value of waste concrete converted to highest-quality RCA \\
$\alpha$ & Sensitivity of waste generators to collection price \\
$\beta$ & Competition intensity between manufacturer and recycler, $\beta \in [0,1]$ \\
$\mu_M$ & Unit transportation cost of manufacturer $M$ \\
$\mu_R$ & Unit transportation cost of recycler $R$ \\
$p_M$ & Unit collection price paid by manufacturer $M$ to waste generators \\
$p_R$ & Unit collection price paid by recycler $R$ to waste generators \\
$q_M$ & Waste concrete collection quantity of manufacturer $M$ \\
$q_R$ & Waste concrete collection quantity of recycler $R$ \\
$n$ & Proportion of R\&D investment cost borne by recycler $R$, $n \in [0,1]$ \\
$k$ & RCA conversion rate of recycler $R$, $k \in [0,1]$ \\
$c$ & R\&D investment cost coefficient \\
$m$ & Actual operational cost of recycler $R$ per unit waste concrete \\
$s$ & Government subsidy per unit waste concrete collected \\
$w$ & Unit outsourcing fee paid by manufacturer $M$ to recycler $R$ \\
$G$ & Total carbon emission quota allocated by government to recycler $R$ \\
$e_0$ & Baseline unit carbon emissions \\
$\gamma$ & Technical emission reduction coefficient \\
$e(k)$ & Dynamic carbon emissions per unit \\
$p_c$ & Trading price per unit of carbon emission in the carbon market \\
$\Pi_i$ & Profit function of participant $i$, $i \in \{M,R\}$ \\
$\varphi_i$ & Shapley value (profit allocation) for participant $i$, $i \in \{M,R\}$ \\
\bottomrule
\end{tabular*}
\begin{tablenotes}[flushleft]
\small\linespread{1}\selectfont
\item \textit{Note}: All decision variables and parameters are non-negative unless otherwise specified.
\end{tablenotes}
\end{threeparttable}
\end{table}
\vspace{10pt}
\normalsize

Based on the above assumptions, the profit functions are expressed as:
\begin{equation}
\Pi_M = (\theta k - p_M - w - \mu_M + s)(\alpha p_M - \beta p_R) - (1 - n)\frac{c k^2}{2}
\end{equation}
\begin{equation}
\Pi_R = (\theta k - p_R - m - \mu_R + s)(\alpha p_R - \beta p_M) + (w - m)(\alpha p_M - \beta p_R) - n\frac{c k^2}{2} - p_c\left[ (e_0 - \gamma k)(\alpha - \beta)(p_M + p_R) - G \right]
\end{equation}

\subsection{Non-cooperative game module}
In this module, there is no R\&D investment sharing between the manufacturer and the recycler. Specifically, the recycler $R$ bears all the R\&D investment costs ($n=1$), while the manufacturer $M$ does not participate in the investment. Under this setting, the recycler independently optimizes the conversion rate $k$ to maximize its own profit, and subsequently both participants compete in collection prices.

Under $n=1$, the profit functions simplify to:
\begin{equation}
\Pi_M^N = (\theta k - p_M - w - \mu_M + s)(\alpha p_M - \beta p_R)
\label{eq:profit_MN}
\end{equation}
\begin{equation}
\Pi_R^N =
\begin{aligned}
& (\theta k - p_R - m - \mu_R + s)(\alpha p_R - \beta p_M) + (w - m)(\alpha p_M - \beta p_R) \\
& - \frac{c k^2}{2} - p_c\left[ (e_0 - \gamma k)(\alpha - \beta)(p_M + p_R) - G \right]
\end{aligned}
\label{eq:profit_RN}
\end{equation}

We solve this game using backward induction. The proof of Proposition 1 is provided in Appendix A. We now present the main result as follows.

\begin{prop}
\label{prop:baseline}
Under the non-cooperative game module (baseline), the equilibrium strategies are given by:
\begin{enumerate}[label={(\arabic*)}]
\item The optimal conversion rate chosen by the recycler is:
\begin{equation}
k^{N*} = \frac{\theta(\alpha p_R^{N*} - \beta p_M^{N*}) + p_c \gamma (\alpha - \beta)(p_M^{N*} + p_R^{N*})}{c}
\label{eq:k_N_star}
\end{equation}

\item The equilibrium collection prices $(p_M^{N*}, p_R^{N*})$ are obtained by Cramer's rule:
\begin{equation}
p_M^{N*} = \frac{N_M^{N}}{D}, \qquad p_R^{N*} = \frac{N_R^{N}}{D},
\label{eq:p_N_star}
\end{equation}
where the numerators and common denominator are
\begin{equation}
\begin{aligned}
N_M^{N} =&\ c^{2}\big[2\alpha^{2}(s - w - \mu_{M}) - \alpha\beta(m + \mu_{R} + e_{0} p_{c} - s) + \beta^{2}(m + e_{0} p_{c} - w)\big] \\
&+ c \cdot \big\{ \alpha^{3}\big[\gamma^{2} p_{c}^{2}(\mu_{M} - s + w) + \gamma p_{c}\theta(2\mu_{M} - \mu_{R} - s + 2w - m) \\
&\qquad - \theta^{2}(m - \mu_{M} + \mu_{R} - w) - e_{0} p_{c}\theta(\gamma p_{c} + \theta)\big] \\
&\quad + \alpha^{2}\beta\big[2\gamma^{2} p_{c}^{2}(-\mu_{M} + s - w) + \gamma p_{c}\theta(3m - 2\mu_{M} + 2\mu_{R} - 3w) \\
&\qquad + \theta^{2}(m - w) + e_{0} p_{c}\theta(3\gamma p_{c} + \theta)\big] \\
&\quad + \alpha\beta^{2}\big[\gamma^{2} p_{c}^{2}(\mu_{M} - s + w) + \gamma p_{c}\theta(-\mu_{R} + s + 2w - 3m) \\
&\qquad - \theta^{2}(m + \mu_{R} - s) - e_{0} p_{c}\theta(3\gamma p_{c} + \theta)\big] \\
&\quad + \beta^{3}\big[\gamma p_{c}\theta(m - w) + \theta^{2}(m - w) + e_{0} p_{c}\theta(\gamma p_{c} + \theta)\big] \big\}, \\[3pt]
N_R^{N} =&\ c^{2}\big[2\alpha^{2}(s - m - \mu_{R} - e_{0} p_{c}) + \alpha\beta(2m - \mu_{M} + s - 3w + 2e_{0} p_{c})\big] \\
&+ c \cdot \big\{ \alpha^{3}\big[2e_{0}\gamma p_{c}^{2}\theta + \gamma^{2} p_{c}^{2}(-\mu_{M} + s - w) + \gamma p_{c}\theta(2m - \mu_{M} + 2\mu_{R} - s - w)\big] \\
&\quad + \alpha^{2}\beta\big[-2e_{0} p_{c}\theta(2\gamma p_{c} + \theta) + 2\gamma^{2} p_{c}^{2}(\mu_{M} - s + w) \\
&\qquad + \gamma p_{c}\theta(2\mu_{M} - 2\mu_{R} - 4m + 4w) - \theta^{2}(2m - \mu_{M} + 2\mu_{R} - s - w)\big] \\
&\quad + \alpha\beta^{2}\big[2e_{0} p_{c}\theta(\gamma p_{c} + \theta) + \gamma^{2} p_{c}^{2}(-\mu_{M} + s - w) \\
&\qquad + \gamma p_{c}\theta(-\mu_{M} + s + 2m - 3w) + 2\theta^{2}(m - w)\big] \big\}, \\[3pt]
D =&\ \alpha^{4}\gamma p_{c}\theta(\gamma p_{c} + \theta)^{2} \\
&+ \alpha^{3}\big[-2\beta\gamma p_{c}\theta(\gamma p_{c} + \theta)^{2} - \beta\theta^{4} - 2c(\gamma p_{c} + \theta)^{2} - 4c\gamma p_{c}\theta - \beta\gamma p_{c}\theta^{3}\big] \\
&+ \alpha^{2}\big[3\beta c(\gamma p_{c} + \theta)^{2} + 4c^{2}\big] + \alpha\beta^{2} c\theta(4\gamma p_{c} + \theta) \\
&+ 2\alpha\beta^{3}\gamma p_{c}\theta(\gamma p_{c} + \theta)^{2} + \alpha\beta^{3}\theta^{4} + \alpha\beta^{3}\gamma p_{c}\theta^{3} \\
&- \beta^{4}\gamma p_{c}\theta(\gamma p_{c} + \theta)^{2} - \beta^{3} c(\gamma p_{c} + \theta)^{2} - \beta^{2} c^{2}.
\end{aligned}
\tag{N1}
\end{equation}

\item The equilibrium collection quantities are:
\begin{equation}
q_M^{N*} = \alpha p_M^{N*} - \beta p_R^{N*}
\quad \text{and} \quad
q_R^{N*} = \alpha p_R^{N*} - \beta p_M^{N*}
\label{eq:q_N_star}
\end{equation}

\item The equilibrium profits are:
\begin{equation}
\Pi_M^{N*} = (\theta k^{N*} - p_M^{N*} - w - \mu_M + s)q_M^{N*}
\label{eq:Pi_MN_star}
\end{equation}
\begin{equation}
\Pi_R^{N*} = (\theta k^{N*} - p_R^{N*} - m - \mu_R + s)q_R^{N*} + (w - m)q_M^{N*}
- \frac{\big[\theta q_R^{N*} + p_c \gamma (\alpha - \beta)(p_M^{N*} + p_R^{N*})\big]^2}{2c} - p_c\left[ (e_0 - \gamma k^{N*})(\alpha - \beta)(p_M^{N*} + p_R^{N*}) - G \right]
\label{eq:Pi_RN_star}
\end{equation}
\end{enumerate}
\end{prop}

 The proof is provided in Appendix A.

Note that the non-cooperative game model characterizes the case where there is no R\&D investment sharing between the concrete manufacturer and the recycler. The recycler independently decides the conversion rate $k$ based on the anticipated collection prices, while both participants engage in price competition for collection. This model serves as a baseline benchmark for subsequent comparative analyses with the biform game model.

\subsection{Two-stage biform game model}
\label{sec:biform_model}
In this subsection, we extend the non-cooperative game module to a non-cooperative and cooperative two-stage biform game model. Building on the scenario in the previous section, we assume the concrete manufacturer is willing to establish a R\&D investment partnership with the recycler. Based on this, we develop a waste concrete recycling management framework based on the biform game, combining non-cooperative and cooperative games. Both participants agree to adopt this biform game structure to formulate their co-opetition relationships.

\subsubsection{Cooperative game}
\label{sec:cooperative_game}

In this subsection, given the pricing strategy profile $(p_M, p_R)$ determined in the non-cooperative stage, we compute the characteristic value $V_{p_M,p_R}(\mathbb{C})$ for any possible coalition $\mathbb{C} \in \{\emptyset, \{M\}, \{R\}, \{M,R\}\}$ using the maximin principle. According to the maximin principle, a coalition $\mathbb{C}$ determines its own optimal decision for maximizing its profit, which is minimized by its opposing coalition. Thus, the characteristic function represents the minimum guaranteed profit a coalition can achieve when facing the worst-case decisions made by the opposing coalition.

For an empty coalition, the profit is certainly zero, thus $V_{p_M,p_R}(\emptyset) = 0$. Next, we calculate $V_{p_M,p_R}(\mathbb{C})$ for $\mathbb{C} \neq \emptyset$. Taking coalition $\{M\}$ as an example, since it acts as a single-player coalition, recycler $R$ forms the opposing coalition. The characteristic value $V_{p_M,p_R}(\{M\})$ depends on $M$'s cost-sharing decision $n$ and $R$'s conversion rate decision $k$, assuming the competitive environment $(p_M, p_R)$ is given. Thus, we apply the maximin principle as follows:
\begin{equation}
V_{p_M,p_R}(\{M\}) = \max_{n \in [0,1]} \min_{k \in [0,1]} \{\Pi_M(p_M, p_R, k, n)\}
\end{equation}
By solving this (detailed in Appendix B), we can obtain the characteristic value $V_{p_M,p_R}(\{M\})$. For other coalitions, we follow the same logic. We present $V_{p_M,p_R}(\mathbb{C})$ for all coalitions in Proposition 2.

\begin{prop}
\label{prop:characteristic}
Under any competitive situation $(p_M, p_R)$, the characteristic values of all possible coalitions $\mathbb{C}$ ($\mathbb{C} \neq \emptyset$) are given as follows:
\begin{equation}
V_{p_M,p_R}(\{M\}) = (s - p_M - w - \mu_M)(\alpha p_M - \beta p_R)
\end{equation}
\begin{equation}
\begin{aligned}
V_{p_M,p_R}(\{R\}) =& (s - p_R - m - \mu_R)(\alpha p_R - \beta p_M) + (w - m)(\alpha p_M - \beta p_R) \\
&+ \frac{\big[\theta(\alpha p_R - \beta p_M) + p_c \gamma (\alpha - \beta)(p_M + p_R)\big]^2}{2c} - p_c e_0 (\alpha - \beta)(p_M + p_R) + p_c G
\end{aligned}
\end{equation}
\begin{equation}
\begin{aligned}
V_{p_M,p_R}(\{M,R\}) =& (s - p_M - m - \mu_M)(\alpha p_M - \beta p_R) + (s - p_R - m - \mu_R)(\alpha p_R - \beta p_M) \\
&+ \frac{(\theta + p_c \gamma)^2(\alpha - \beta)^2(p_M + p_R)^2}{2c} - p_c e_0 (\alpha - \beta)(p_M + p_R) + p_c G
\end{aligned}
\end{equation}
\end{prop}

 The proof is provided in Appendix B.

The coalitions' characteristic values computed in Proposition 2 constitute the cooperative game for R\&D investment cooperation. In cooperative game theory, the characteristic value reflects a coalition's bargaining power. To apply the Shapley value to determine a fair profit allocation scheme, it is crucial to check whether this cooperative game is convex and super-additive. If so, the core of this game is non-empty, and the Shapley value, as a single solution, will fall into the core.

\begin{prop}
\label{prop:convexity}
The cooperative game defined by the characteristic functions in Proposition 2 is convex and super-additive. Specifically, for any $(p_M, p_R)$, we have:
\begin{equation}
\begin{aligned}
V_{p_M,p_R}(\{M,R\}) - V_{p_M,p_R}(\{M\}) - V_{p_M,p_R}(\{R\}) =&\ \frac{\theta^2}{2c}(\alpha p_M - \beta p_R)[(\alpha p_M - \beta p_R) + 2(\alpha p_R - \beta p_M)] \\
&+ \frac{p_c \gamma}{c} (\alpha - \beta)(p_M + p_R) \big[\theta(\alpha - \beta)(p_M + p_R) - \theta(\alpha p_R - \beta p_M)\big]
\end{aligned}
\end{equation}
Consequently, $V_{p_M,p_R}(\{M,R\}) \geq V_{p_M,p_R}(\{M\}) + V_{p_M,p_R}(\{R\})$, and the Shapley value allocation belongs to the core.
\end{prop}

 The proof is provided in Appendix C.

Proposition 3 implies that forming the grand coalition generates more profit, providing a strong incentive for both the concrete manufacturer and the recycler to cooperate. 
The Shapley value is an important and widely accepted allocation mechanism in the cooperative game theory. Although cooperative games with a nonempty core offer several distribution strategies, this paper utilizes the Shapley value due to its superior practical acceptance among participants. This allocation method is effective because it aligns individual rewards with each player's distinct contribution to the overall coalition. The Shapley value allocations $\varphi_M$ and $\varphi_R$ are calculated as:

\begin{equation}
\begin{aligned}
\varphi_M(p_M, p_R) =& (s - p_M - m - \mu_M)(\alpha p_M - \beta p_R) \\
&+ \frac{\theta^2(\alpha p_M - \beta p_R)\big[(\alpha p_M - \beta p_R) + 2(\alpha p_R - \beta p_M)\big]}{4c} \\
&+ \frac{\theta p_c \gamma (\alpha - \beta)(p_M + p_R)(\alpha p_M - \beta p_R)}{2c}
\end{aligned}
\end{equation}
\begin{equation}
\begin{aligned}
\varphi_R(p_M, p_R) =& (s - p_R - m - \mu_R)(\alpha p_R - \beta p_M) + (w - m)(\alpha p_M - \beta p_R) \\
&+ \frac{\theta^2(\alpha p_R - \beta p_M)^2}{4c} + \frac{\theta^2(\alpha - \beta)^2(p_M + p_R)^2}{4c} \\
&+ \frac{\theta p_c \gamma (\alpha - \beta)(p_M + p_R)\big[(\alpha p_M - \beta p_R) + 2(\alpha p_R - \beta p_M)\big]}{2c} \\
&+ \frac{p_c^2 \gamma^2 (\alpha - \beta)^2 (p_M + p_R)^2}{2c} - p_c e_0 (\alpha - \beta)(p_M + p_R) + p_c G
\end{aligned}
\end{equation}

In this sense, the players can achieve optimal profits in the cooperative game part, making it feasible to deploy the non-cooperative game by using these resulting allocations as payoff functions.

\subsubsection{Non-cooperative game}
\label{sec:noncooperative_game}

Following the biform game logic, the resulting allocations derived in the cooperative game part give the payoff functions in the non-cooperative game part. With the competitive environment defined by the pricing strategies, players anticipate how cooperation profit will be allocated. Therefore, the payoff functions for $M$ and $R$ in the non-cooperative price competition stage are given by:
\begin{equation}
f_M(p_M, p_R) = \varphi_M(p_M, p_R)
\end{equation}
\begin{equation}
f_R(p_M, p_R) = \varphi_R(p_M, p_R)
\end{equation}

By using the Nash equilibrium concept, we maximize these two payoff functions. We take the first-order partial derivatives of $f_M(p_M, p_R)$ and $f_R(p_M, p_R)$ with respect to $p_M$ and $p_R$, respectively, and set them equal to zero. Solving these reaction functions (detailed in Appendix D) yields the equilibrium solutions as shown in Proposition 4.

\begin{prop}
\label{prop:biform}
Under the two-stage biform game model, the equilibrium collection prices $(p_M^{C*}, p_R^{C*})$ are obtained by Cramer's rule:
\begin{equation}
p_M^{C*} = \frac{N_M^{C}}{D^{C}}, \qquad p_R^{C*} = \frac{N_R^{C}}{D^{C}},
\end{equation}
where
\begin{equation}
\begin{aligned}
N_M^{C} =&\ c^{2}\big[8\alpha^{2}(s - m - \mu_{M}) - 4\alpha\beta(m + \mu_{R} + e_{0} p_{c} - s) + 4\beta^{2}(m + e_{0} p_{c} - w)\big] \\
&+ c \cdot \big\{ \alpha^{3}\big[8\gamma^{2} p_{c}^{2}(m + \mu_{M} - s) + 2\gamma p_{c}\theta(3m + 4\mu_{M} - \mu_{R} - 3s) \\
&\qquad + 2\theta^{2}(m + 2\mu_{M} - \mu_{R} - s) - 2e_{0} p_{c}\theta(\gamma p_{c} + \theta)\big] \\
&\quad + \alpha^{2}\beta\big[16\gamma^{2} p_{c}^{2}(-m - \mu_{M} + s) + 2\gamma p_{c}\theta(-3m - 6\mu_{M} + 2\mu_{R} + 4s - w) \\
&\qquad + 2\theta^{2}(-2\mu_{M} + \mu_{R} + s - w) + 2e_{0} p_{c}\theta(3\gamma p_{c} + 2\theta)\big] \\
&\quad + \alpha\beta^{2}\big[8\gamma^{2} p_{c}^{2}(m + \mu_{M} - s) + 2\gamma p_{c}\theta(-m + 2\mu_{M} - \mu_{R} - s + 2w) \\
&\qquad + 2\theta^{2}(-m + \mu_{M} - \mu_{R} + w) - 2e_{0} p_{c}\theta(3\gamma p_{c} + 2\theta)\big] \\
&\quad + \beta^{3}\big[2\gamma p_{c}\theta(m - w) + 2\theta^{2}(m - w) + 2e_{0} p_{c}\theta(\gamma p_{c} + \theta)\big] \big\}, \\[3pt]
N_R^{C} =&\ c^{2}\big[8\alpha^{2}(s - m - \mu_{R} - e_{0} p_{c}) + 4\alpha\beta(m - \mu_{M} + s - 2w + 2e_{0} p_{c})\big] \\
&+ c \cdot \big\{ \alpha^{3}\big[4e_{0}\gamma p_{c}^{2}\theta + 2e_{0} p_{c}\theta^{2} + 8\gamma^{2} p_{c}^{2}(-m - \mu_{M} + s) \\
&\qquad + 2\gamma p_{c}\theta(-m - 3\mu_{M} + 2\mu_{R} + s) + 2\theta^{2}(-\mu_{M} + \mu_{R})\big] \\
&\quad + \alpha^{2}\beta\big[-2e_{0} p_{c}\theta(4\gamma p_{c} + 3\theta) + 16\gamma^{2} p_{c}^{2}(m + \mu_{M} - s) \\
&\qquad + 4\gamma p_{c}\theta(2m + 6\mu_{M} - 2\mu_{R} - 4s + 2w) + 2\theta^{2}(3\mu_{M} - 2\mu_{R} - s + w)\big] \\
&\quad + \alpha\beta^{2}\big[4e_{0} p_{c}\theta(\gamma p_{c} + \theta) + 8\gamma^{2} p_{c}^{2}(-m - \mu_{M} + s) \\
&\qquad + 2\gamma p_{c}\theta(-m - 3\mu_{M} + 3s - 2w) + 2\theta^{2}(m - \mu_{M} + s - 2w)\big] \big\}, \\[3pt]
D^{C} =&\ \alpha^{4}\big[2\gamma^{3} p_{c}^{3}\theta + 5\gamma^{2} p_{c}^{2}\theta^{2} + 4\gamma p_{c}\theta^{3} + \theta^{4}\big] \\
&+ \alpha^{3}\beta\big[-4\gamma^{3} p_{c}^{3}\theta - 10\gamma^{2} p_{c}^{2}\theta^{2} - 8\gamma p_{c}\theta^{3} - 2\theta^{4}\big] \\
&+ \alpha^{3} c\big[-8\gamma^{2} p_{c}^{2} - 24\gamma p_{c}\theta - 12\theta^{2}\big] + \alpha^{2}\beta c\big[12\gamma^{2} p_{c}^{2} + 24\gamma p_{c}\theta + 12\theta^{2}\big] + 16\alpha^{2} c^{2} \\
&+ \alpha\beta^{3}\big[4\gamma^{3} p_{c}^{3}\theta + 10\gamma^{2} p_{c}^{2}\theta^{2} + 8\gamma p_{c}\theta^{3} + 2\theta^{4}\big] + \alpha\beta^{2} c\theta(8\gamma p_{c} + 4\theta) \\
&- \beta^{4}\big[2\gamma^{3} p_{c}^{3}\theta + 5\gamma^{2} p_{c}^{2}\theta^{2} + 4\gamma p_{c}\theta^{3} + \theta^{4}\big] - \beta^{3} c(4\gamma^{2} p_{c}^{2} + 8\gamma p_{c}\theta + 4\theta^{2}) - 4\beta^{2} c^{2}.
\end{aligned}
\tag{C1}
\end{equation}
Substituting $(p_M^{C*}, p_R^{C*})$ back into the grand coalition's optimal decisions, the optimal conversion rate $k^{C*}$ and R\&D investment sharing proportion $n^{C*}$ are:
\begin{equation}
k^{C*} = \frac{(\theta + p_c \gamma)(\alpha-\beta)(p_M^{C*} + p_R^{C*})}{c}
\end{equation}
\begin{equation}
n^{C*} = 1 - \frac{\theta^{2} q_M^{C*}\big(3 q_M^{C*} + 2 q_R^{C*}\big)}{2(\theta + p_c \gamma)^{2} \big[(\alpha-\beta)(p_M^{C*} + p_R^{C*})\big]^{2}} - \frac{\theta p_c \gamma \, q_M^{C*}}{(\theta + p_c \gamma)^{2} (\alpha-\beta)(p_M^{C*} + p_R^{C*})}
\end{equation}
where $q_{M}^{C*} = \alpha p_M^{C*} - \beta p_R^{C*}$ and $q_{R}^{C*} = \alpha p_R^{C*} - \beta p_M^{C*}$. Equivalently, the optimal $n^{C*}$ is uniquely determined by equating each participant's actual cooperative profit under the chosen sharing rule to the corresponding Shapley allocation. Specifically, substituting the cooperative profits and $k^{C*}$ into $\Pi_M^{actual}(n) = \varphi_M$ yields
\begin{equation}
n^{C*} = 1 - \frac{2c \big[(\theta k^{C*} - p_M^{C*} - w - \mu_M + s) q_M^{C*} - \varphi_M(p_M^{C*}, p_R^{C*})\big]}{(\theta + p_c \gamma)^2 \big[(\alpha - \beta)(p_M^{C*} + p_R^{C*})\big]^2},
\end{equation}
which can be verified to coincide with the analogous expression obtained from $\Pi_R^{actual}(n) = \varphi_R$ (a verification given in Appendix D). When $n^{C*}$ lies outside $[0,1]$, the corner solutions $n^{C*} = 0$ or $n^{C*} = 1$ apply, reflecting parameter regions where the Shapley allocation cannot be sustained by an interior sharing rule.
\end{prop}

 The proof is provided in Appendix D.

By employing the Shapley value, we find a fair allocation solution which suggests reasonable profit functions for the non-cooperative game. The biform game enables us to derive analytical results for a complex problem where collection pricing competition and R\&D investment cooperation interact—a dynamic that is almost impossible to solve comprehensively using traditional single-stage game approaches.

\section{Numerical simulation}
To illustrate the equilibrium outcomes more intuitively and investigate the impacts of key parameters (including price sensitivity coefficient, competition intensity, recycling value, and transportation costs) on collection prices, collection quantities, cost-sharing ratios, and technology improvement levels, we conduct numerical simulations using MATLAB 2023b.

\subsection{Parameter settings}
Waste concrete recycling represents a global challenge, and parameter values vary across different countries and contexts. To ensure the numerical simulation aligns with practical scenarios, we situate our analysis in the Chinese context. This selection is justified by two considerations: First, China has demonstrated strong policy commitment to waste concrete recycling, exemplified by initiatives such as the ``14th Five-Year Plan'' for solid waste pollution prevention and control, which explicitly mandates construction waste resource utilization rates to reach 60\% in large and medium-sized cities by 2025. Second, China's waste concrete recycling industry is at an early stage of development, with firms actively pursuing technological upgrades and seeking R\&D collaborations, which parallels the non-cooperative and cooperative game framework in this study. Accordingly, this context provides valuable insights for firm decision-making and government policy guidance.

Based on China's official policies, industry reports, project disclosures, expert interviews, and relevant literature, we calibrate the parameters as presented in Table~\ref{tab:parameters}.

\begin{table}[!htbp]
\centering
\captionsetup{font=normalsize, labelsep=period}
\caption{Parameter settings and baseline values}
\label{tab:parameters}
\normalsize
\begin{tabular}{lp{10cm}c}
\toprule
\textbf{Parameter} & \textbf{Source} & \textbf{Value} \\
\midrule
$\alpha$ & Drawing upon the Survey Report on the Resource Utilization of Construction Waste in Shanghai, the collection price for waste concrete typically ranges from 10 to 15 yuan per ton. Following the parameter specification in \citep{wang2025Mechanism}'s approach, we set $\alpha$ = 15/10 = 1.5 & 1.5 \\
\midrule
$\beta$ & According to the China Construction Waste Recycling Industry Overview, the compound annual growth rate of construction waste is approximately 13.4\%. Given the substantial capital requirements and technical barriers, the number of major firms remains relatively stable. We interpret the annual incremental waste concrete resources as the competition intensity. & 0.134 \\
\midrule
$m$ & Based on the report in World Environment journal, processing 100,000 cubic meters of waste concrete costs approximately 8 million yuan. We convert the unit. & 32 yuan/m$^3$ \\
\midrule
$w$ & Drawing from annual reports of A-share listed firms specializing in waste concrete recycling (e.g., Sinoma Resources, China Tianying, Shanghai Environment), the gross profit margins range from 7.88\% to 13.04\%. Taking the average margin of 9.84\%, we calculate $w = m \times (1 + 9.84\%) = 35$ yuan/m$^3$. & 35 yuan/m$^3$ \\
\midrule
$s$ & According to the 2025 project disclosure for the Fuzhou (Jiangxi Province) construction waste treatment franchise, the government provides a subsidy of 35 yuan per ton based on collection volume. We convert the unit. & 87.5 yuan/m$^3$ \\
\midrule
$\mu_M$ & We conducted market research by requesting quotes from multiple logistics providers and interviewing construction cost professionals. & 25 yuan/m$^3$ \\
\midrule
$\mu_R$ & We conducted market research by requesting quotes from multiple logistics providers and interviewing construction cost professionals. & 44 yuan/m$^3$ \\
\midrule
$c$ & Based on equipment pricing data from the China Construction Waste Recycling Industry Overview (including crushers, vibrating feeders, and magnetic separators), we calculated the weighted average cost coefficient. & 135,000 \\
\midrule
$\theta$ & According to Liu's survey report, the selling price of recycled concrete aggregate from waste concrete ranges from 50 to 60 yuan per ton. Taking the median value and converting units. & 137.5 yuan/m$^3$ \\
\midrule
$p_c$ & Referring to the carbon trading prices in various Chinese pilot markets in 2021, which ranged from 50 to 80 yuan per ton. Taking the median value. & 65 yuan/ton \\
\midrule
$G$ & We follow the parameter specification in the research of \cite{sun2021Optimal}. & 300 ton \\
\midrule
$e_0$ & According to the European Environment Agency (EEA), processing 1 ton of waste concrete generates approximately 0.02 to 0.05 tons of carbon dioxide emissions. Taking the median value as baseline. & 0.035 ton CO$_2$/m$^3$ \\
\midrule
$\gamma$ & Since direct empirical data for this specific technical coefficient is difficult to obtain, we follow the parameter setting approach for numerical simulations adopted by \citep{he2026EmissionReduction} & 0.015 \\
\bottomrule
\end{tabular}
\end{table}
\vspace{-15pt}

\subsection{Simulation results and analysis}
\label{sec:simulation_analysis}

To investigate the impacts of key parameters on equilibrium outcomes, we conduct a comprehensive numerical simulation analysis. Based on the baseline parameter values specified in Section 4.1, we vary each parameter within a reasonable range while holding other parameters constant. The simulation results are presented in Figs. 1-5.

\subsubsection{Comparison between Non-cooperative and Cooperative Games}
\label{sec:comparison}

This subsection presents the numerical comparison between the non-cooperative game model and the biform game model under varying key parameters. Specifically, we evaluate whether enterprises should choose to cooperate in a competitive environment from two perspectives: economic profit and environment and resource conversion.

Figure 1 evaluates the enterprises' decisions from the perspective of economic profit. Figure 1(a) illustrates the effects of price sensitivity ($\alpha$) and competition intensity ($\beta$) on optimal profits. As price sensitivity increases, both enterprises experience an upward profit trajectory across both game models. While heightened sensitivity typically suppresses equilibrium prices, it simultaneously triggers a substantial expansion in total market demand, compensating for reduced unit margins. Notably, the cooperative strategy consistently yields a Pareto improvement, and the profit gap between the biform and non-cooperative models widens as sensitivity grows. This occurs because joint R\&D investment lowers marginal processing costs, enabling the allied firms to maintain robust margins. Conversely, escalating competition intensity ($\beta$) exerts downward pressure on the profits of both participants. This reflects a classic prisoner's dilemma in the front-end collection market, where fierce competition inevitably shrinks margins. However, the cooperative alliance prevents destructive price wars and provides a robust buffer against these competitive frictions, sustaining a relatively stable cooperative surplus. Establishing a strategic alliance is vital for the recycler to mitigate risks and ensure economic viability.

The analysis demonstrates that forming a cooperative alliance allows enterprises to leverage price-sensitive demand expansion, effectively avoiding the prisoner's dilemma in collection markets \citep{lin2025Cooperation}. Furthermore, by linking surplus distribution directly to each firm's operational contribution, the Shapley value mechanism highlights the recycler's specialized processing capacity as the primary driver of market value, which extends the fairness allocation principles proposed by \cite{jia2023Novel}. By sharing R\&D investments, the cooperative mechanism inherently prevents destructive price wars and ensures that the recycler maintains stable profit margins even as front-end competition intensifies. Ultimately, this strategic partnership enables the recycler to survive severe competitive pressures while remaining an independent processing enterprise.

Figure 1(b) illustrates the influence of the outsourcing fee ($w$) and the government subsidy ($s$) on the economic profit of the two enterprises. Regarding the outsourcing fee ($w$), when the fee is relatively low, both participants choose to cooperate because the biform game yields a higher surplus than the non-cooperative game. However, as the outsourcing fee increases and exceeds a specific threshold, the manufacturer's profit under the cooperative alliance falls below its non-cooperative baseline, prompting the manufacturer to abandon cooperation. Regarding the government subsidy ($s$), for the manufacturer, there are two crossing points: when the subsidy is too low, the cooperative returns are insufficient to compensate the manufacturer for the shared R\&D investments, leading to a non-cooperative preference. At moderate subsidy levels, the manufacturer strongly prefers the cooperative alliance. However, when the subsidy becomes excessively high, the manufacturer's independent non-cooperative profit grows faster than its Shapley value allocation, again making the non-cooperative strategy more lucrative. In contrast, the recycler only faces a single threshold, transitioning from non-cooperation to cooperation once the subsidy reaches a minimum level and consistently benefiting thereafter. Thus, to ensure both parties choose to cooperate, the subsidy must be maintained within a moderate optimal interval.

Identifying these threshold boundaries challenges the conventional assumption that higher external subsidies unconditionally benefit recycling networks \citep{yuan2023Differentiated}. By simultaneously modeling front-end pricing competition and back-end R\&D investment cooperation, this analysis defines the boundary conditions for alliance stability. It reveals that excessive external subsidies crowd out the manufacturer's incentive to share R\&D investments, as their independent profits outpace cooperative gains. Additionally, the analysis demonstrates that an excessive internal outsourcing fee disrupts the delicate profit-sharing equilibrium between collection and treatment, with the manufacturer bearing the initial brunt of the destabilization. This mechanism explains the fragility and instability of joint innovation alliances often observed in third-party remanufacturing \citep{chen2025Exploring}. Ultimately, these findings indicate that securing a sustainable cooperative recycling ecosystem requires carefully calibrated inter-firm pricing contracts and moderate, targeted policy interventions rather than indiscriminate financial support.

Figure 1(c) illustrates the impact of the potential recycled market value ($\theta$) and the transportation cost asymmetry ($\Delta\mu$) on the economic performance of both enterprises. Regarding the market value ($\theta$), both the concrete manufacturer and the recycler consistently prefer the cooperative alliance across the entire range. An escalation in $\theta$ expands the total economic surplus, and the biform game enables both participants to capture a larger portion of this wealth. This occurs because shared R\&D investments upgrade processing technologies, allowing the alliance to capture the surging market value more efficiently than operating independently. Regarding the transportation cost asymmetry ($\Delta\mu$), when the cost asymmetry is small, both participants experience a cooperative deficit and prefer the non-cooperative strategy. As the asymmetry increases to a moderate level, both participants secure higher profits under the cooperative strategy, establishing a stable "sweet spot" for the alliance. However, if the cost gap widens to an extreme level, the manufacturer continues to benefit, but the recycler's cooperative profit drops below its non-cooperative baseline, causing the recycler to exit the partnership. Moderate asymmetry balances the manufacturer's reverse logistics advantage with the recycler's technical specialization, whereas extreme asymmetry skews the profit distribution and destabilizes the cooperative equilibrium.

By incorporating the construction industry's logistical asymmetries into the payoff functions, this analysis reveals how cost differences shape cooperative stability. The identification of a moderate cost-asymmetry "sweet spot" refines facility integration strategies \citep{liu2019Expand} by demonstrating that extreme logistical differences can actually disrupt coopetition \citep{lin2025Cooperation}. Furthermore, while previous literature notes a positive relationship between end-product valuation and recycling viability \citep{li2020Key}, this biform game approach specifies the underlying mechanism: higher market values inherently motivate voluntary R\&D cooperation. Consequently, policy-driven market enhancements \citep{su2024Stakeholder}, such as green public procurement protocols or standardized certification systems, can effectively shift the financial burden of technological upgrades from government mandates to shared, market-driven investments.

Figure 2 evaluates the strategic decisions from the perspective of environment and resource conversion by analyzing the impacts of key parameters on the conversion rate ($k$). Figure 2(a) illustrates the effects of the potential recycled market value ($\theta$) and the technical emission reduction coefficient ($\gamma$) on the conversion rates. As the potential market value ($\theta$) increases, the conversion rates under both models rise. The cooperative conversion rate grows at a faster pace than its non-cooperative counterpart, which expands the absolute difference between the two models. A higher end-product value provides the economic incentive for shared R\&D investments, leading to higher processing efficiency. Conversely, an increase in the technical emission reduction coefficient ($\gamma$) improves the conversion rates in both scenarios by lowering the carbon compliance burden. Under this condition, the non-cooperative rate grows slightly faster than the cooperative rate, causing the ratio between the two models to decrease. When the technology offers higher emission reduction efficiency, the recycler faces reduced environmental compliance pressure, which enables them to improve their independent R\&D capabilities. Nevertheless, the conversion rate of the cooperative model remains higher than the non-cooperative baseline.

Figure 2(b) illustrates the effects of price sensitivity ($\alpha$) and competition intensity ($\beta$) on the conversion rates. An increase in price sensitivity ($\alpha$) expands market demand, which drives up the conversion rates for both models. Because an active market incentivizes independent entities to increase their processing capabilities, the non-cooperative conversion rate grows slightly faster, leading to a decrease in the ratio between the two models. Regarding the competition intensity ($\beta$), as front-end collection competition intensifies, the conversion rates in both models decline. However, the non-cooperative model experiences a more significant decrease. Independent entities under high competitive pressure tend to reduce their independent R\&D budgets to maintain profitability. In contrast, the cooperative alliance shares the R\&D investment, which buffers against the front-end market competition and results in a smaller decline in conversion efficiency. Consequently, the relative advantage of the cooperative model increases under high competition. This indicates that in competitive markets, the cooperative mechanism helps protect R\&D investments and sustain higher waste conversion rates.

In summary, the concrete manufacturer and the recycler in the waste concrete recycling process make distinct strategic choices depending on their primary evaluation perspectives. From the perspective of economic profitability, establishing a cooperative alliance is not always the most effective decision. Influenced by operational and external factors such as government subsidies, outsourcing fees, and transportation cost asymmetries, the enterprises may face cooperative deficits and rationally choose the non-cooperative strategy to protect their margins. However, from the perspective of environmental performance and resource conversion, cooperation is consistently the optimal choice. If the objective is to achieve higher waste concrete conversion rates and reduce treatment-related carbon emissions through technological R\&D, the enterprises will strictly prefer to cooperate.

\subsubsection{Impact of Cap-and-Trade Mechanism}
\label{sec:cap_and_trade}

This section evaluates how the cap-and-trade mechanism impacts the decisions of the waste manufacturer and the recycler within the waste concrete recycling process. Figures 3(a) and 3(b) illustrate the impacts of the carbon trading price and the actual carbon emissions per unit on the optimal profits of the concrete manufacturer and the recycler, respectively. As both the carbon trading price and the unit carbon emissions escalate, the economic returns for both participants exhibit a downward trajectory. Notably, the profit of the recycler is more sensitive to the escalating carbon trading price, experiencing a much steeper and more profound decline compared to the manufacturer. This reflects the recycler's direct operational exposure and heavy financial burden regarding emission compliance. While an increase in unit carbon emissions also erodes profits across the board, its overall negative impact is milder than that of the carbon price, though it still affects the specialized recycler more distinctly than the manufacturer.

The comparative analysis of the two game reveals the advantage of the cooperative alliance. Under the biform game model, both the magnitude and the rate of profit decline for the two enterprises are smaller than those observed in the non-cooperative game. This reveals that R\&D investment cooperation effectively acts as a robust financial buffer. By establishing a cooperative alliance, the concrete manufacturer and the recycler can mitigate the negative financial shocks induced by volatile carbon prices and stringent emission penalties, thereby preserving their profit margins and enhancing the economic resilience of the entire waste concrete recycling network.

Fig. 4 demonstrates how the participants react when the carbon trading price increases. As the price goes up, the recycler faces a much higher cost to process the waste concrete because they are the actual party generating emissions. To balance this extra expense, the recycler has to reduce the collection price offered to the market. The concrete manufacturer does not bear the emission cost, so their collection price stays stable. The collection quantities also drop and it pulls down the total quantities. Furthermore, the heavier financial burden leaves the recycler with less capital to share the R\&D investment. This leads to a lower conversion rate. Ultimately, the rising carbon price cuts into the profit margins of both enterprises, but the recycler suffers a much sharper decline due to their direct exposure to the carbon market.

These findings align with prior research showing that carbon pricing mechanisms alter the cost structure of the supply chains. Under the cap-and-trade mechanism, enterprises directly exposed to emission compliance costs bear disproportionate burdens, creating structural inequalities in inter-firm cooperation \citep{bai2022Distributionally}. The differential vulnerability between manufacturers and recyclers mirrors observations in EV battery closed-loop supply chains where carbon costs reshape pricing dynamics across recycling channels \citep{wu2024Electric}. Our results extend this insight by demonstrating that rising carbon prices compress margins and undermine the financial prerequisites for cooperative R\&D, reducing both the conversion rate and the R\&D investment sharing proportion. This suggests that in high carbon-price environments, policy interventions may be necessary to maintain cooperative incentives in recycling networks \citep{huang2024Optimal}.

Fig. 5 illustrates the impact of the baseline unit carbon emissions $e_0$ on the optimal profits. A higher $e_0$ means the recycler starts with more carbon emissions per unit before technology improvements. This inflates the total carbon compliance cost in the same way a higher trading price does. Facing higher operational costs, the recycler scales back their front end procurement price. The shrinking market volume and increased cost pressure make both participants less willing to achieve R\&D investment cooperation, causing a steady decline in the conversion rate and the sharing ratio. The profit trajectories confirm that if the processing technology is not environmentally friendly, the resulting carbon costs will damage the economic profit of cooperation in the waste recycling process.

This finding agrees with previous studies on emission reduction in waste management. When the processing technology generates a large amount of carbon emissions, the recycler faces high compliance costs. These high costs leave the recycler with less money for R\&D \citep{besklubova2026Establishing}. Because processing waste concrete becomes very expensive under high emission levels, both the concrete manufacturer and the recycler lose the economic incentive to invest in R\&D. Consequently, the conversion rate and the investment sharing proportion paid by the recycler both decrease. The results also show that the R\&D investment cooperation is easily broken by high carbon costs. The enterprises can easily lower their collection prices to survive sudden cost increases, but they will quickly stop their joint R\&D investments when profit margins shrink. This proves that developing processing technologies with low carbon emissions is strictly necessary to keep the cooperative partnership profitable and stable \citep{chen2025Lowcarbon}.

In summary, the cap-and-trade mechanism introduces an environmental cost that reshapes the cost structure of the concrete manufacturer and the recycler. Because the recycler processes the waste concrete, they are the most vulnerable to high carbon prices and poor emission efficiencies. When these carbon costs become too heavy, the recycler is forced to shrink their collection quantities and reduce R\&D investments, which negatively impacts the entire waste concrete recycling supply chain.

\section{Conclusion and limitations}
\subsection{Conclusions}

This study develops a two-stage biform game model to investigate the coopetition relationships between the concrete manufacturer and the recycler in the waste concrete recycling process under the cap-and-trade mechanism. By comparing the equilibrium solutions of the biform game model with a non-cooperative game model, the study draws the following principal conclusions:

First, whether the concrete manufacturer and the recycler choose to cooperate in a competitive environment depends fundamentally on whether they evaluate the decision from the perspective of economic profit or environment and resource conversion. From the perspective of economic profit, establishing a cooperative alliance is not always the most effective decision. Influenced by operational and external factors such as government subsidies, outsourcing fees, and transportation cost asymmetries, the enterprises may choose the non-cooperative strategy to protect their margins. However, from the perspective of environment and resource conversion, cooperation is consistently the optimal choice. If the objective is to achieve higher waste concrete conversion rates and reduce treatment-related carbon emissions through technological R\&D, the enterprises will prefer to cooperate.

Second, concerning the impacts of the cap-and-trade mechanism, the study finds that an increase in the carbon trading price generally raises costs and reduces profit margins of the concrete manufacturer and the recycler. Cooperation in the competition environment provides a buffering effect against such volatility. By sharing the R\&D investment, the manufacturer enables the recycler to reach a higher conversion rate, which generates greater marginal revenue and enhances the whole supply chain's resilience to rising carbon market costs.

Theoretically, this study constructs a two-stage biform game model for waste concrete recycling, providing a novel analytical method and theoretical perspective for the field of resource recycling and management. Practically, the findings encourage resource recycling participants to establish cooperative relationships during the recycling process and incorporate internal carbon risk-sharing mechanisms. Furthermore, policymakers are advised to implement targeted structural subsidies, achieving a sustainable balance between environmental protection and the incentivization of enterprise economic activities.

\subsection{Limitations and future research directions}
There are certain limitations that highlight avenues for future research. First, this study primarily focuses on a basic dual-channel supply chain consisting of a single traditional manufacturer and a single specialized recycler. However, the actual waste concrete recycling ecosystem often involves more complex interactions among multiple stakeholders. Future research could expand the current modeling framework to incorporate multiple competing recyclers, third-party logistics providers, or explicitly consider the downstream demand market's acceptance of recycled materials, thereby more comprehensively revealing the operational mechanisms of a closed-loop supply chain network with multi-party participation.

Second, this research employs a static game framework to analyze the equilibrium decisions and profit distribution of supply chain members. Future research could integrate differential games or evolutionary game theory to explore the dynamic evolutionary trajectories of enterprises' technology R\&D investments over multiple periods, as well as the long-term stability conditions of the outsourcing cooperation.

\section*{Data availability statement}
All data, models, or codes that support the findings of this study are available from the corresponding author upon reasonable request.

\section*{Acknowledgments}
This paper was supported by the National Natural Science Foundation of China (Project Numbers: 72371189, 72371190) and the program of China Scholarship Council (CSC) (No. 202206260227).

\section*{Supplemental Material}
Appendix are available online at https://doi.org/10.5281/zenodo.20528901

\bibliographystyle{elsarticle-harv}
\bibliography{references}


\end{document}
